% Source-derived chapter 08: Further issues
% Original source: very-stable-higgs-bundles-equivariant-multiplicity-mirror-symmetry
\chapter{Further issues}
\label{very-stable-higgs-bundles-equivariant-multiplicity-mirror-symmetry-chapter-08}
\source{very-stable-higgs-bundles-equivariant-multiplicity-mirror-symmetry}

\begin{remark}[Source section]
\label{very-stable-higgs-bundles-equivariant-multiplicity-mirror-symmetry-chapter-08-source}
\source{very-stable-higgs-bundles-equivariant-multiplicity-mirror-symmetry}
This chapter records the source section; no Lean declaration has yet been checked for this material.
\notready
\end{remark}

% The source section title was: Further issues
\label{further}
\source{very-stable-higgs-bundles-equivariant-multiplicity-mirror-symmetry}


\subsection{Simple Lie groups}\label{simple}
\source{very-stable-higgs-bundles-equivariant-multiplicity-mirror-symmetry}
The simplest example of a very stable upward flow of type $(1,1,...1)$ is, as we have seen, the canonical section of the Hitchin fibration. In fact  the moduli space of Higgs bundles for any simple Lie group $\G$ admits   a corresponding section, constructed in \cite{hitchin-lie}. For the adjoint group this is unique. This points towards a generalization of the $(1,1,...1)$ type in this context, where the mirror should be a hyperholomorphic bundle on the moduli space for the Langlands dual group $\G^{\vee}$.

The key point in \cite{hitchin-lie} is to use the principal three-dimensional subgroup of $\G$ \cite{Kos}. The fixed point $\mathcal E$ is then obtained from the canonical uniformising $\SL_2$ Higgs bundle  via the homomorphism $\SL_2\rightarrow \G$. At a fixed point  the scalar action on the Higgs field arises from the action of a subgroup of the adjoint  group of  $\G$ and in this case it is the  action of the image of $\C^\times\subset \SL_2$.

The characteristic feature of the subgroup is that it breaks up the Lie algebra $\lie{g}$ into $\ell$ irreducible components where $\ell$ is the rank of $\G$ and the dimensions are $2d_i-1$ where the $d_i$ are the degrees of a basis of invariant polynomials. These are odd-dimensional and hence representations of $\SO_3$. Therefore $t   \in \T$ acts through $(t^{-1},1,t)\subset \SO_3$.

The  Higgs field  has homogeneity one with respect to the $\T$-action. Each  irreducible representation of $\SO_3$  has a single weight one subspace so the Higgs field must take values  in an $\ell$-dimensional subspace of $\lie{g}$. From \cite[Lemma 5.2]{Kos}, the action has  weight $1$ on the root spaces of the $\ell$ simple roots $\alpha_1,\dots, \alpha_{\ell}$ which therefore form a basis for this subspace.  In the case of the general linear group  treated in this paper   the simple roots are $x_{i+1}-x_i$ which correspond to line bundles  $L_{i+1}L_i^*$ and we dealt with  sections $b_i$ of $L_{i+1}L_i^*K$. The direct analogue for the group $\G$  of the $(1,1,\dots,1)$ type  fixed point is then a line bundle $L_{\alpha_i}$ for each simple root and a section $b_i$ of $L_{\alpha_i}K$. Writing any root $\alpha$ in terms of simple roots then defines by tensor product a line bundle $L_{\alpha}$ and a Lie algebra bundle
$$H\oplus \bigoplus_{\alpha\in \Delta}L_{\alpha}$$
where $H$ is a trivial rank $\ell$ bundle associated to the Cartan subalgebra. For the canonical fixed point $\calE_0$  in the moduli space of $\G$-bundles each $b_i$ is nonvanishing and the upward flow is a section of the Hitchin fibration, and in particular is closed. We shall discuss next the case of a Higgs bundle $\calE$ when $b_i$ is of degree $m_i\ge 0$ by studying the virtual multiplicity $$m_\calE(t)=\frac{\chi_{\T}(\Sym ({T_\calE^+}^*))}{\chi_{\T}(\Sym ({\mathcal A}^*))}.$$
Here $T_\calE^+$ is the sum of root space line bundles for roots of positive weight under the $\C^{\times}$-action. Since the weight is $1$ on the simple roots, these are the positive roots. If  $\alpha=c_1\alpha_1+\dots+c_{\ell}\alpha_{\ell}$ in terms of the simple roots then $w(\alpha)=c_1+\dots+c_{\ell}$.   Since $m_i=\deg L_{\alpha_i}K$ we have
$$\deg L_{\alpha}=\sum_{i=1}^{\ell} c_i(\alpha) (m_i-(2g-2)). $$


Following the procedure of  Remark 5.12 we obtain
$$\chi_\T(T_\calE^+)=\sum_{\alpha>0} t^{-w(\alpha)}\left(\sum_{w(\alpha)=k}-(\deg L_{\alpha}+(1-g))+\sum_{w(\alpha)=k-1}\deg(L_{\alpha}+(g-1))\right)$$
or
$$\chi_\T(T_\calE^+)=\sum_{\alpha>0} t^{-w(\alpha)}\left(\sum_{w(\alpha)=k}\sum_{i=1}^{\ell} -m_ic_i(\alpha) +\sum_{w(\alpha)=k-1}\sum_{i=1}^{\ell} m_ic_i(\alpha) \right)+(g-1)N(t).$$
When all the $m_i=0$, the multiplicity is $1$ because then we have an equivariant section of ${\mathcal A}\rightarrow {\mathcal M}$ (one may also easily do the calculation -- see below the generating function for $N_j$) so the $(g-1)N(t)$ term above cancels the denominator ${\chi_{\T}(\Sym ({\mathcal A}^*))}$ and the multiplicity reduces to
\begin{equation}
m_\calE(t)=\prod_{i=1}^{\ell}\prod_{\alpha>0}\left(\frac{1-t^{w(\alpha)+1}}{1-t^{w(\alpha)}}\right)^{m_ic_i(\alpha)}.
\label{multG}
\source{very-stable-higgs-bundles-equivariant-multiplicity-mirror-symmetry}
\end{equation}
By experimentation one finds that this is rarely a polynomial. In the case of ${\G}_2$ this happens independently of the $m_i$, supposing they are not identically zero, for other groups there are constraints on the $m_i$. So there is a distinct difference from the case of $\GL_n$ or $\SL_n$. One case however always gives a polynomial as we see next.

A simple root $\alpha_i$ is called {\it cominuscule} if the coefficient $c_i(\alpha)=0,\pm1$ for every root $\alpha$.  Every simple root of $\SL_n$ has this property but for other types  they are special or non-existent. The coroot $\alpha_i^{\vee}$ in the dual root system, the root system of the Langlands dual group, is called {\it minuscule} (see for example \cite{Buch}) and this corresponds to an irreducible representation: the exterior powers of the vector representation for $\SL_n$, the vector representation for $\mathrm{Sp}_{2n}$, the spin representation for $\mathrm{Spin}_{2n+1}$ and the vector and two spin representations for $\mathrm{Spin}_{2n}$.


Suppose then that $\alpha_i$ is cominuscule and the Higgs field has $m_j=0$ for $j\ne i$ and $m_i=1$. The formula (\ref{multG}) is multiplicative in $m_i$ so this case is sufficient. Since $c_i(\alpha)=0,1$ for a positive root, the multiplicity is
$$\prod_{\alpha>0, c_i(\alpha)\ne 0}\left(\frac{1-t^{w(\alpha)+1}}{1-t^{w(\alpha)}}\right). $$
Consider this product over all the positive roots of $\G$:
$$\prod_{\alpha>0}\left(\frac{1-t^{w(\alpha)+1}}{1-t^{w(\alpha)}}\right). $$
Let $N_j$ be the number of root spaces of weight $j$, then this expression is
$$q_\G(t)=\prod_{j\ge 1}\left(\frac{1-t^{j+1}}{1-t^{j}}\right)^{N_j}= \prod_{j\ge 1}(1-t^{j})^{N_{j-1}-N_j}.$$

The Lie algebra $\lie{g}$ splits under the $\SO_3$-action into $\ell$  irreducibles of dimension $2d_i-1$ where $d_i$ is the degree of a basis element of the ring of invariant polynomials. In each irreducible of dimension $2d_i-1$, there is a one-dimensional space of weight $1,2,\dots d_i-1$. Then $$N_1x+N_2x^2+\cdots =\sum_{i=1}^{\ell}(x+x^2+\cdots + x^{d_i-1})=\frac{x}{(1-x)}\sum_1^{\ell}(1-x^{d_i-1}).$$
So the generating function for $N_{j-1}-N_j$ is
$-\ell x+\sum_1^{\ell}x^{d_i}$, hence
\begin{equation}
q_\G(t)=(1-t)^{-{\ell}}  \prod_{j=1}^{\ell}(1-t^{d_j}).
\label{qG}
\source{very-stable-higgs-bundles-equivariant-multiplicity-mirror-symmetry}
\end{equation}

Associated to each cominuscule root is a maximal parabolic subgroup $\rm P$ whose Levi subgroup $\rm L$  is obtained by deleting $\alpha_i$ from the Dynkin diagram of $\G$ to produce a semisimple group and  has an additional $\C^\times$ factor corresponding to the coroot ${\alpha_i}^{\vee}$ (see \cite{Buch} for example). All simple roots of type $A_{\ell}$ are minuscule and we have dealt with those in the main body of the paper. For  the other groups the deleted Dynkin diagram is connected and so we get a simple group. The roots for which $c_i(\alpha)=0$ are the roots of $\rm L$.  The $\T$-action has weight one on  the simple root spaces  of $\G$ and hence also  on those of $\rm L$. This means that $\T$ acts with weight one on a linear combination $\sum_{j\ne i} \lambda_je_{\alpha_j}$ of root vectors. From \cite{Kos} if $\lambda_j\ne 0$, this is a principal nilpotent in the Levi subalgebra and so $\T$ acts as a semisimple element in a principal three-dimensional subgroup of the Levi group $\rm L$.

Let $n_j$ be the degrees of  generators of the invariant polynomials on $\rm L$. The rank of $\rm L$ is $\ell$ and the simple component $\ell-1$.  Then
$$q_{\rm L}(t)=(1-t)^{-{\ell}}  \prod_{j=1}^{\ell-1}(1-t^{n_j})$$
and
\begin{equation}
m_\calE(t)=\prod_{\alpha>0, c_i(\alpha)\ne 0}\left(\frac{1-t^{w(\alpha)+1}}{1-t^{w(\alpha)}}\right)=\frac{q_\G(t)}{q_{\rm L}(t)}= \prod_{j=1}^{\ell}\frac{(1-t^{d_j})}{(1-t^{n_j})}.
\label{mminusc}
\source{very-stable-higgs-bundles-equivariant-multiplicity-mirror-symmetry}
\end{equation}
\begin{proposition}
\source{very-stable-higgs-bundles-equivariant-multiplicity-mirror-symmetry}
\notready \label{cominint} The virtual multiplicity $m_\calE(t)$ in equation (\ref{mminusc}) is a polynomial.
\end{proposition}
\begin{proof} From \cite[Theorem 26.1]{Borel}  if $K,U$ are compact Lie groups of the same rank then the Poincar\'e polynomial of $K/U$ is
	$$P(K/U,t)=\frac{(1-t^{s_1})(1-t^{s_2})\cdots (1-t^{s_{\ell}})}{(1-t^{r_1})(1-t^{r_2})\cdots (1-t^{r_{\ell}})}.$$
	where $s_i-1, r_i-1$ are the degrees of generators of the cohomology of $K,U$. Taking $K$ and $U$ to be the maximal compact subgroups of ${\rm P},{\rm L}$ then $s_j=2d_j, r_j=2n_j$ and   $P(K/U,t)$ is a polynomial $Q(t^2)$. Then $m_\calE(t)=Q(t)$.
\end{proof}
A theorem (which also follows from the geometric Satake theorem of \cite{mirkovic-villonen}, c.f. \cite[Theorem 1.6.3]{ginzburg}) of B.Gross \cite[Corollary 6.4]{Gross}  states that the Lefschetz action of $\SL_2$ on the cohomology of the flag variety $\G/{\rm P}$ is isomorphic to the representation of the principal $\SL_2$ of the Langlands dual $\G^{\vee}$ on the  representation given by the minuscule weight dual to $\alpha_i$. In the light of our results on exterior powers for the linear group this suggests the following.

\begin{conjecture}
\source{very-stable-higgs-bundles-equivariant-multiplicity-mirror-symmetry}
\notready  \label{cominmirror} Let $\alpha_{i_1},\dots,\alpha_{i_k}$ be distinct cominuscule roots and suppose that the corresponding sections $b_{i_j}$ of $L_{\alpha_{i_j}}K$ have $m$ distinct zeros $x_1,\dots, x_m$, and for other simple roots $b_i$ has no zero.
	Then the fixed point  is very stable and the mirror of its upward flow is the tensor product of the vector bundles associated to the universal principal bundle for the Langlands dual group at $x_j$ in the corresponding minuscule representation. \end{conjecture}

	\begin{remark}
\source{very-stable-higgs-bundles-equivariant-multiplicity-mirror-symmetry}
\notready
The referee has observed a further example beyond the general linear group supporting this conjecture. The cominuscule root for $\mathrm{Sp}_{2n}$ gives a fixed point which is a rank $2n$ symplectic Higgs bundle of the form
$$L_1\stackrel{1}\rightarrow L_2\stackrel{1}\rightarrow\cdots \stackrel{1}\rightarrow L_n\stackrel{b_n}\rightarrow L_n^{-1}\stackrel{1}\rightarrow \cdots \stackrel{1}\rightarrow L_2^{-1}\stackrel{1}\rightarrow L_1^{-1}$$
	and if $b_n$ has no multiple zeros this is very stable as an $\mathrm {SL}_{2n}$-bundle and a fortiori as an $\mathrm{Sp}_{2n}$-bundle.

	On the upward flow from this point, the Higgs field $\Phi$
	at a zero of $b_n$  preserves a Lagrangian subspace. At a regular point the eigenspaces of $\Phi$ lie in pairs  in $n$ $2$-dimensional symplectic subspaces so there are $2^n$ invariant Lagrangian subspaces which should correspond  to a basis for the spin representation of the Langlands dual group $\mathrm{Spin}_{2n+1}.$

\end{remark}


\subsection{Multiple zeros} \label{multiple}
\source{very-stable-higgs-bundles-equivariant-multiplicity-mirror-symmetry}

We showed in Theorem 1.2 that a necessary and sufficient condition for a fixed point of type $(1,1,\dots,1)$ to be very stable is that $b$
should have no multiple zeros. It is natural to ask what happens when we do have a multiple zero. What is the closure of the upward flow,
and is there a corresponding hyperholomorphic bundle on the mirror moduli space?

We already know from Remark 5.15 that the intersection
of the upward flow with a generic fibre has fewer points than the situation where $b$ has simple zeros. We also see in the proof of Lemma 4.13
that a multiple zero leads to a Hecke curve in the closure -- the closure of the flow tangent to a nilpotent Higgs field.

We illustrate this  here with the simplest
example of this situation in rank $2$, for the group $\SL_2$. Let ${\mathcal E}$ be a fixed point of the $\T$-action defined by the vector bundle
$L\oplus L^*$ and the off-diagonal Higgs
field $b: L\rightarrow L^*K$, defining  a section of $L^{-2}K$ which we assume has  a single double zero at $c\in C$, and hence $L\cong K^{1/2}(-c)$.
The upward
flow is given
by Higgs bundles $(E,\Phi)$ where $E$ is an extension $L\rightarrow E\stackrel p\rightarrow L^*$, $\tr \Phi=0$ and $b=p \Phi:L\rightarrow L^*K$.

The spectral curve  $ C_a$ is defined by $\det(x-\Phi)=x^2+\det \Phi=0$ and $\sigma(x)= -x$ is the covering involution. Then if $ C_a$ is smooth, as we have seen, $E=\pi_*U$ where the line
bundle $U$ is the cokernel of $x-\Phi:\pi^*(E\otimes K^*)\rightarrow \pi^*E$. The inclusion $L\subset E\in H^0(C;\Hom(L,E))$ is canonically defined
by the direct image construction as a section of the line bundle $\Hom (\pi^*L,U)$ on $C_a$.

Now $b$ vanishes at  $c\in C$ and suppose further that $\det \Phi(c)\ne 0$. Since  $L$ is preserved by $\Phi$ at $c$, it  maps to zero in the cokernel
of $\Phi-x$ for one of the two eigenvalues $\pm x$. Then $p=(x,c)\in  C_a$ or $\sigma(p)=(-x,c)$ are zeros of the holomorphic section of $(\pi^*L)^*U$.
Since $\pi:C_a\rightarrow C$ is unramified at these points, these zeros have multiplicity $2$ since $b$ has a double zero at $c$.


The bundle $\Lambda^2 E$ is trivial which implies $U\cong M\pi^*K^{1/2}$ where $\sigma^*M\cong M^*$ (i.e. $M$ lies in the Prym variety) and so
since $L\cong K^{1/2}(-c)$, $\pi^*L^* U\cong M(c)$ which has degree $2$. But we have a section which vanishes with multiplicity $2$
hence $M(c)\sim 2p$ or $2\sigma(p)$ as divisors, or, since $c\sim p+\sigma(p)$, $M\sim \pm(p-\sigma(p))$. Thus the
intersection with the generic fibre consists of two points in the Prym variety whereas if $b$ had two distinct zeros the intersection
would be four points.

Note that if $q(c)=0$ and $C_a$ is smooth, there is a single point $p$ with $\pi(p)=c$ and $\pi$ locally has the form $z\mapsto z^2$. Since $b$ has a double zero at $c$, the section of $\pi^*L^* U$ has a simple zero which is a contradiction since it has degree $2$. Thus the upward flow does not intersect these fibres, but as we shall see next, the closure does.

Suppose $(E,\Phi)$ is a limit point of the upward flow. There is still a non-zero homomorphism $K^{1/2}(-c)\rightarrow E$ but its image may not be an
embedding. It will generate a subbundle $L_0$ of larger degree. Since it has positive degree by  stability it cannot be invariant by the Higgs
field $\Phi_0$ so there is a nonvanishing section of $L_0^{-2}K$. But $\deg L_0>\deg K^{1/2}(-c)=g-2$ hence $\deg L_0^{-2}K\le 0$.
The existence of a section means that we have equality and $L_0=K^{1/2}$, so the bundle is an extension $K^{1/2}\rightarrow E\rightarrow
K^{-1/2}$ and $b$ is an isomorphism. But from \cite[Proposition 3.3]{hitchin} the extension is trivial and the limit lies on the canonical section. Since
$\det \Phi\in  H^0(C,K^2)$ parametrizes the section the limit lies on
the $(3g-4)$-dimensional  linear subspace consisting of sections which vanish at $c$.


Consider now the uniformising Higgs bundle $K^{1/2}\oplus K^{-1/2}$ with Higgs field $1: K^{1/2}\rightarrow K^{-1/2}K$. The {\it downward}
flow from this point consists of extensions $K^{-1/2}\rightarrow E\rightarrow K^{1/2}$ with the same nilpotent Higgs field mapping the quotient to the subbundle. This is a Lagrangian submanifold and the
extensions are parametrized by a class $e\in H^1(C,K^*)$ (of  dimension $(3g-3)$ = $\dim {\mathcal M}/2$). The section $s$
of ${\mathcal O}(c)$
vanishing at $c$ defines a homomorphism $s:K^{1/2}(-c)\rightarrow K^{1/2}$ which lifts to $E$ if $se=0\in  H^1(C,K^*)$. Equivalently
$e=\delta(u)$ for
$u\in H^0(c,{\mathcal O}_c(K^*))\cong \C$ where $\delta$ is the connecting homomorphism in the cohomology exact sequence.

In Dolbeault terms with a $C^{\infty}$ splitting $K^{1/2}\oplus K^{-1/2}$ we have
$$\bar\partial+ \begin{pmatrix}0 & \bar\partial v/s\cr
0 & 0\end{pmatrix}\qquad \Phi= \begin{pmatrix}0 & 1\cr
0 & 0\end{pmatrix} $$
where (as in 4.1.4) $v$ is a local extension of $u$, zero outside a neighbourhood and holomorphic in a smaller one.  Then the embedding of $K^{1/2}(-c)\subset E$ is $(v,-s)$ so that $\Phi(-v,s)=(s,0)$ and
the skew form
on $E$ gives
$\langle (v,-s),(s,0)\rangle= s^2$. This is equivalent to projecting onto the quotient and clearly has a double zero at $c$. Then each extension $E$
lies on the upward flow from ${\mathcal E}$ and has as limit the canonical fixed point. This is a projective line in the closure, a Hecke curve.


We see therefore that the union of the upward flow $W^+_c$ from $c$ and $W^+_0$ from the canonical point is a closed subspace, and the symplectic form vanishes at  smooth points.

\begin{conjecture}
\source{very-stable-higgs-bundles-equivariant-multiplicity-mirror-symmetry}
\notready  The mirror
	of the closed Lagrangian cycle $W^+_0\cup W^+_c$ is the universal adjoint bundle for $c$ in the $\SO_3$ moduli space. \end{conjecture}

Here is some motivation.  The subvariety $W^+_0\cup W^+_c$ intersects a generic fibre in three points and so the Fourier-Mukai transform is generically a rank $3$
vector bundle. The virtual multiplicity $m_\calE(t)$ for a double zero of $b$ is the same as for two distinct zeros and, from the point of view of the previous subsection this is  $(1+t)^2$, from the Poincar\'e polynomial of $\PP^1\times \PP^1$. Moreover the Lefschetz action of $\SL_2$ on the cohomology is the action on $V\otimes V$ where $V$ is the defining 2-dimensional representation. This is the direct sum of the adjoint  and the trivial one-dimensional
representation so on the adjoint component the action is consistent with the conjecture.

Now the  adjoint bundle for $\SO_3$ is a rank $3$ bundle $V$ and $\Phi: K^*\rightarrow V$  is injective when $\Phi$ is regular at
each point of $C$, i.e. when the spectral curve is smooth. Trivializing $K$
at $c$, the universal bundle admits a homomorphism from the trivial bundle. It follows that the canonical section, which corresponds to the trivial line bundle on the spectral curve,  forms a component of the
support of the inverse transform. The other component clearly relates to the quotient sheaf which is generically a rank $2$ bundle.

One point to note is that the hyperholomorphic connection on the universal bundle is not  reducible (its Pontryagin class is a generator of $H^4({\mathcal M},\R)$) and this casts a question mark over the
issue of whether the mirror of the closure of $W^+_c$ actually has a hyperholomorphic connection.

More generally, there is a Hecke correspondence for every irreducible representation of the Langlands dual group, which has a natural compactification (see \cite[\S 9]{kapustin-witten}).  In our case, the image of the Hitchin section in the $\SL_2$ Higgs moduli space under this compactified Hecke  correspondence, attached to the adjoint representation of $\SO_3$, is precisely our Lagrangian cycle $W_0^+\cup W_c^+$. Its mirror then should be the universal adjoint bundle at $c$. The $\T$-equivariant index of its structure sheaf then should be $1+t+t^2$, which is also the $\T$-character of the mirror adjoint bundle at the canonical uniformising Higgs bundle. Incidently, it also matches the intersection Poincar\'e polynomial of a certain weighted projective space, which is the closure of the space of Hecke modifications at a point in our case (c.f. \cite[9.2]{kapustin-witten}). In fact, from the geometric Satake correspondence \cite{mirkovic-villonen,ginzburg} this generalises to the general case of simple (even reductive) groups, and we expect the analoguous results to hold there too.

\subsection{The cotangent fibre} \label{cotfib}
\source{very-stable-higgs-bundles-equivariant-multiplicity-mirror-symmetry}

The upward flow from a very stable bundle $(E,0)$ is an important example.
We have calculated the rank of its transform in Remark~\ref{multN} using the multiplicity formula. The alternative method, as in \cite{beauville-etal} is to argue
that if there are no nilpotent Higgs fields then locally the moduli space is given by the vanishing of invariant polynomials. For a simple Lie group this is therefore the vanishing of  $(2d_i-1)(g-1)$ functions of degree $d_i$ for $i=1,\dots,\ell$ and so the rank of this bundle is
$$\prod_1^{\ell}d_i^{(2d_i-1)(g-1)},$$  which Laumon in \cite[Remark 6.2.3]{laumon2} called the global analogue of the order of the Weyl group.
The structure of this bundle even for the $2^{3g-3}$ case of $SL_2$ is largely unknown, and less so a way to construct a hyperholomorphic connection on it.  In \eqref{cotmir}  we offered evidence from mirror symmetry for a conjecture of the $\T$-character of the  fibre of this bundle with a $\T$-equivariant structure at the fixed points of type $(1,1)$.

It is an important example as  it relates to the work of Donagi--Pantev \cite{DP} on the
geometric Langlands programme, where they proposed a geometric approach of constructing automorphic sheaves on the moduli stack of bundles, which was, in fact, also the motivation of  Drinfeld \cite{drinfeld} and Laumon \cite{laumon} for the introduction of very stable bundles, see also \cite[Conjecture 6.2.1]{laumon2}. This programme aims to associate to a point  $m\in {\mathcal M}^{\vee}$, the mirror of the Higgs bundle moduli space ${\mathcal M}$, a flat connection on ${\mathcal N}$, the moduli space of stable bundles, with singularities on the ``wobbly locus". A generic point in the mirror consists of a spectral curve $C_a$ and (according to the SYZ procedure) a line bundle $L$ over its Jacobian. In ${\mathcal M}$, the spectral curve defines a fibre of the Hitchin fibration. This fibre over $a\in {\mathcal A}$ intersects $T^*{\mathcal N}\subset {\mathcal M}$ in the complement of a subvariety and the projection to ${\mathcal N}$ extends to a rational map as in \cite{beauville-etal}. The idea is that the vector bundle $V$ on ${\mathcal N}$ is associated to a direct image of $L$, and the canonical
1-form on the cotangent bundle $T^*{\mathcal N}$ defines a Higgs field on $V$, a section of $\End (V)\otimes T^*$. Then, using the higher dimensional nonlinear Hodge theory due to Simpson and  Mochizuki \cite{mochizuki}, one hopes to obtain a flat connection.

What is the link with this paper? The complement of the wobbly locus in ${\mathcal N}$ is the set of very stable vector bundles $E$. The fibre of $V$  at a generic $[E]\in {\mathcal N}$ is the direct sum of the fibres of $L$ over the intersection of the upward flow of $(E,0)$ with the Hitchin fibre. Put another way, if ${\mathcal N}^{vs}\subset {\mathcal N}$ denotes the moduli space of very stable bundles, then the Fourier-Mukai transform of the upward flows over  ${\mathcal N}^{vs}$
defines a bundle on ${\mathcal N}^{vs}\times {\mathcal M}^{\vee}$. In this paper we have been concerned with the restriction to $[E]\times {\mathcal M}^{\vee}$, for Donagi and Pantev it is the restriction to ${\mathcal N}^{vs}\times \{m\}$.

In fact what we have encountered is a miniature version of this in the case of very stable Higgs bundles of type $(1,1,\dots, 1)$. Instead of the $\T$-fixed point set ${\mathcal N}$ we have another component: a subvariety defined by the zero set of $(b_1,\dots, b_{n-1})$ in the product of symmetric products $C^{[m_1]}\times C^{[m_2]}\times \cdots \times C^{[m_{n-1}]}$. Moreover we have explicitly  identified  the wobbly locus, namely where the product $b_1\dots b_{n-1}$ has multiple zeroes. This  includes the various diagonals. The hyperholomorphic bundle we identified as a tensor product of exterior powers of universal bundles corresponding to the zeros of the $b_i$. By nonlinear Hodge theory on $C$, a point $m\in {\mathcal M}^{\vee}$ defines a flat connection on the curve $C$ and hence also on the product $C^{m_1}\times C^{m_2}\times \cdots \times C^{m_{n-1}}$, invariant by the action of the product of symmetric groups, hence descending to a flat connection, but with singularities over the diagonals.
